% !TEX root = task-2-1-2.tex
% Editable task bank. Append new BankTask blocks at the end of the document.
\documentclass[12pt,a4paper]{article}
\usepackage{physics-textbook}
\usepackage{graphicx,needspace}
\usetikzlibrary{decorations.pathmorphing}
\newif\ifShowSolutions
\ifdefined\TutorEdition
  \ShowSolutionstrue
\else
  \ShowSolutionsfalse % Change to \ShowSolutionstrue to display solutions.
\fi
\newcounter{banktask}
\newsavebox{\BankTaskBox}
\newcommand{\BankTaskContent}[3]{%
  \noindent\textbf{\thebanktask.}\quad #2\par
  \ifShowSolutions
    \Needspace{50mm}% Keep the solution heading with its opening text or diagram.
    #3\par
    \ifdefined\TutorEdition
      \ifcsname TutorSource#1\endcsname
        \nopagebreak\smallskip{\small\csname TutorSource#1\endcsname\par}%
      \fi
    \fi
  \fi
}
\newcommand{\BankTask}[3]{%
  \refstepcounter{banktask}\par\medskip
  \sbox{\BankTaskBox}{\begin{minipage}{\linewidth}%
    \BankTaskContent{#1}{#2}{#3}\end{minipage}}%
  \ifdim\dimexpr\ht\BankTaskBox+\dp\BankTaskBox\relax>\dimexpr\textheight-20mm\relax
    \Needspace{35mm}\BankTaskContent{#1}{#2}{#3}%
  \else
    \Needspace{\dimexpr\ht\BankTaskBox+\dp\BankTaskBox\relax}%
    \noindent\usebox{\BankTaskBox}\par
  \fi
  \medskip
}
\begin{document}
\BookSubsection{2.1.2}{Тепловое движение атомов и молекул}


% Source sdamgia-805: https://phys-ege.sdamgia.ru/problem?id=805
\BankTask{sdamgia-805}{
Какое из утверждений правильно?

$A$. Диффузия наблюдается только в газах и жидкостях.

$B$. Диффузия наблюдается только в твердых телах.

$C$. Диффузия наблюдается в газах, жидкостях и твердых телах.

$1$) $A$

$2$) $B$

$3$) $C$

$4$) ни $A$, ни $B$, ни $C$
}{
\smallskip\noindent\textbf{Решение}\par
Диффузия — это процесс взаимного проникновения молекул одного вещества между молекулами другого за счет теплового движения. Тепловое движение молекул присуще всем агрегатным состояниям, так что диффузия наблюдается и в газах, и в жидкостях, и в твердых телах.

Ответ: $3$.
}

% Source sdamgia-3286: https://phys-ege.sdamgia.ru/problem?id=3286
\BankTask{sdamgia-3286}{
Броуновское движение частиц пыльцы в воде объясняется

$1$) хаотичностью химических реакций на поверхности частиц

$2$) непрерывностью и хаотичностью теплового движения молекул воды

$3$) существованием сил притяжения и отталкивания между атомами в молекулах

$4$) наличием питательных веществ в воде
}{
\smallskip\noindent\textbf{Решение}\par
Броуновское движение происходит из-за того, что все жидкости и газы состоят из атомов или молекул — мельчайших частиц, которые находятся в постоянном хаотическом тепловом движении, и потому непрерывно толкают броуновскую частицу с разных сторон. Правильный ответ $2$.

Ответ: $2$.
}

% Source sdamgia-3287: https://phys-ege.sdamgia.ru/problem?id=3287
\BankTask{sdamgia-3287}{
Явление диффузии в жидкостях объясняется тем, что молекулы жидкостей

$1$) отталкиваются друг от друга

$2$) колеблются около своих положений равновесия

$3$) притягиваются друг к другу

$4$) могут хаотично перемещаться по объему
}{
\smallskip\noindent\textbf{Решение}\par
Диффузией называется процесс взаимного проникновения молекул одного вещества между молекулами другого, приводящий к самопроизвольному выравниванию их концентраций по всему занимаемому объему. Молекулы жидкости могут хаотично перемещаться по объему (они ведут «кочевой образ жизни»). Именно этим и объясняется главным образом явление диффузии в жидкостях.

Ответ: $4$.
}

% Source sdamgia-3296: https://phys-ege.sdamgia.ru/problem?id=3296
\BankTask{sdamgia-3296}{
Броуновским движением называется

$1$) упорядоченное движение слоев жидкости (или газа)

$2$) упорядоченное движение твердых частиц вещества, взвешенных в жидкости (или газе)

$3$) конвекционное движение слоев жидкости при ее нагревании

$4$) хаотическое движение твердых частиц вещества, взвешенных в жидкости (или газе)
}{
\smallskip\noindent\textbf{Решение}\par
Броуновским движением называется беспорядочное движение микроскопических частиц твердого вещества (пылинки, крупинки взвеси, частички пыльцы растения и так далее), вызываемое тепловым движением частиц жидкости или газа. Верно определение $4$.

Ответ: $4$.
}

% Source sdamgia-3560: https://phys-ege.sdamgia.ru/problem?id=3560
\BankTask{sdamgia-3560}{
Укажите пару веществ, скорость диффузии которых наибольшая при прочих равных условиях:

$1$) раствор медного купороса и вода

$2$) пары эфира и воздух

$3$) свинцовая и медная пластнины

$4$) вода и спирт
}{
\smallskip\noindent\textbf{Решение}\par
Диффузией называется процесс взаимного проникновения молекул одного вещества между молекулами другого за счет непрерывного теплового движения. Поскольку наиболее свободно перемещаться в пространстве могут именно молекулы газов, диффузия газов имеет максимальную скорость. Таким образом, при прочих равных условиях, скорость диффузии наибольшая для паров эфира и воздуха.

Ответ: $2$.
}

% Source sdamgia-4485: https://phys-ege.sdamgia.ru/problem?id=4485
\BankTask{sdamgia-4485}{
Диффузия в жидкости происходит быстрее при повышении температуры, потому что с повышением температуры

$1$) увеличиваются силы взаимодействия молекул

$2$) увеличивается скорость теплового движения молекул

$3$) жидкости расширяются

$4$) уменьшаются силы взаимодействия молекул
}{
\smallskip\noindent\textbf{Решение}\par
Диффузия — это процесс взаимного проникновения молекул одного вещества между молекулами другого, приводящий к самопроизвольному выравниванию их концентраций по всему занимаемому объему. Он происходит за счет непрерывного хаотического движения молекул. Как известно, увеличение температуры приводит к увеличению скорости теплового движения. Поэтому правильным является ответ под номером $2$.

Ответ: $2$.
}

% Source sdamgia-7621: https://phys-ege.sdamgia.ru/problem?id=7621
\BankTask{sdamgia-7621}{
Скорость диффузии зависит от

А) температуры

Б) агрегатного состояния веществ

Из представленных утверждений являются верными

$1$) только А

$2$) только Б

$3$) и А, и Б

$4$) ни А, ни Б
}{
\smallskip\noindent\textbf{Решение}\par
Скорость диффузии зависит и от температуры и от агрегатного состояния вещества.

Ответ: $3$
}

% Source sdamgia-4555: https://phys-ege.sdamgia.ru/problem?id=4555
\BankTask{sdamgia-4555}{
Какое(-ие) из утверждений правильно(-ы)?

$A$. Скорость диффузии увеличивается с ростом температуры.

Б. Скорость движения броуновской частицы не зависит от ее массы.

$B$. Диффузия наблюдается в газах, жидкостях и твердых телах.

$1$) только А

$2$) только Б

$3$) А и В

$4$) Б и В
}{
\smallskip\noindent\textbf{Решение}\par
Диффузия — это процесс взаимного проникновения молекул одного вещества между молекулами другого, приводящий к самопроизвольному выравниванию их концентраций по всему занимаемому объему. Он происходит за счет непрерывного хаотического движения молекул. Как известно, увеличение температуры приводит к увеличению скорости теплового движения: $E= \frac{3}{2} kT.$ Утверждение А верно.

Поскольку тепловое движение молекул свойственно для всех агрегатных состояний, диффузия имеет место и в тверды, и в жидких, и в газообразных тела. Отличается только скорость диффузии. Утверждение В верно.

Броуновское движение — беспорядочное движение микроскопических видимых, взвешенных в жидкости или газе частиц твердого вещества, вызываемое тепловым движением частиц жидкости или газа. Броуновские частицы находятся в тепловом равновесии со средой, а значит их средняя кинетическая энергия движения пропорциональна температуре (работает та же формула, что и выше). Отсюда делаем вывод, что скорость движения броуновской частицы зависит от ее массы, так как $E= \frac{m\overline v ^{2}}{2} .$ Чем больше масса частицы, тем меньше ее скорость. Утверждение Б ошибочно.

Ответ: $3$
}

% Source sdamgia-12857: https://phys-ege.sdamgia.ru/problem?id=12857
\BankTask{sdamgia-12857}{
Броуновская частица массой $1{,}3\cdot 10^{-15}\,\text{кг}$ находится в жидкости при температуре $300\,\text{К}$. Чему равна среднеквадратичная скорость этой частицы, если в системе установилось термодинамическое равновесие? (В мм/с; округлите до целого числа.)
}{
\smallskip\noindent\textbf{Решение}\par
Среднеквадратичная скорость частицы равна:

$v = \sqrt{\frac{3kT}{m_0}} = \sqrt{\frac{3 \cdot 1{,}38 \cdot 10 ^{- 23} \cdot 300}{1{,}3 \cdot 10 ^{- 15}}} \approx 3 \cdot 10 ^{- 3} \text{м}/\text{с}=3\text{мм}/\text{с}.$

Ответ: $3$.
}

% Source sdamgia-23295: https://phys-ege.sdamgia.ru/problem?id=23295
\BankTask{sdamgia-23295}{
\begin{center}\begin{tikzpicture}[task graph]
\draw[task grid,xstep=1.5,ystep=.7] (0,0) grid (4.5,4.9);\draw[task axis,->] (-.2,0)--(4.8,0) node[right] {$t,\,{}^\circ\text{C}$};\draw[task axis,->] (0,0)--(0,5.1) node[above left] {$V^2,\,\text{м}^2/\text{с}^2$};
\node[below left] at (0,0) {$0$};\foreach \x/\s in {0/77,1.5/427,3/777}{\fill (\x,0) circle (1.2pt);}\node[below right] at (0,-.05) {$77$};\foreach \x/\s in {1.5/427,3/777}{\node[below] at (\x,-.05) {$\s$};}\foreach \y/\s in {1.4/311625,2.8/623250,4.2/934875}{\node[left] at (0,\y) {$\s$};}\draw[task curve] (0,1.4)--(3,4.2);\foreach \x/\y in {0/1.4,1.5/2.8,3/4.2}{\fill (\x,\y) circle (1.7pt);}
\end{tikzpicture}\end{center}
На рисунке изображен график зависимости величины среднего значения квадрата скорости молекул идеального газа от температуры. Определите молярную массу этого газа. (В граммах на моль; округлите до целого числа.)
}{
\smallskip\noindent\textbf{Решение}\par
Средний квадрат скорости молекулы равен $\vec v ^{2} = \frac{3RT}{M} .$ Откуда, используя данные из графика, находим: $M= \frac{3RT}{\vec v ^{2}} = \frac{3 \cdot 8{,}31 \cdot 700}{623250} =0{,}028\text{кг}/\text{моль} = 28\text{г}/\text{моль}.$

Ответ: $28$.
}

\end{document}
